% Chapter 2, Figure 12: critically damped response in yaw (scan: figures/ch2/fig12.png).
% Curve: fig12.py, eq. (22) alpha_X = (A_1 + A_2 t) e^{-Dt} with eq. (23) and D = C_2/2I_L = omega_n (zeta = 1);
% alpha_X0 = 1, Omega_X0 = 0.6 in the family's units (time 1/omega_n). The straight line is the true tangent at
% t = 0, slope Omega_X0 (standing rule 4: the 1973 line is steeper than the curve it leaves).
% Layout shared by Figs 10-14 (tamr sketch; 0.4 in per 1/omega_n, 0.576 in per unit of yaw angle;
% the axis below zero shortened from the printed -2.25 to -1.8 units).
\documentclass[11pt]{standalone}
\usepackage{tamrfig}
\pgfplotsset{free response/.style={tamr sketch, x=0.4in, y=0.576in,
    xmin=0, xmax=9.6, ymin=-1.8, ymax=2.8,
    xlabel={$t$ (sec)}, ylabel={$\alpha_X$ (rad)},
    % the label above the arrow, upright (tamr sketch's rotate=-90 turns it on its side under axis lines=middle)
    every axis y label/.style={at={(ticklabel* cs:1.0)}, anchor=south, inner sep=3pt}}}
\tikzset{tangent/.style={s2, line width=1pt}}
\begin{document}
\begin{tikzpicture}
  \begin{axis}[free response, ytick={1}, yticklabels={$\alpha_{X0}$}]
    \addplot[series1] table[x=t, y=alpha, col sep=comma] {figures/v2/ch2/fig12.csv};
    % tangent at t = 0: alpha_X0 + Omega_X0 t
    \draw[tangent] (axis cs:0,1) -- (axis cs:2.3,2.38)
      node[anchor=west, text=ink, font=\small, inner sep=2pt] {Slope${}=\Omega_{X0}$};
    \node[anchor=east, font=\footnotesize, text=ink2, inner sep=3pt] at (axis cs:0,0) {0};
  \end{axis}
\end{tikzpicture}
\end{document}
